\section{Final Wrapping SNARK}\label{sec:wrapping_snark_design}

This section refines the high-level description from Subsection
\ref{subsec:wrapping_snark} into a concrete final proving layer. After the IVC
chain terminates, the prover holds a final folded accumulator, its witness, the
last folding proof, and the data required for deferred memory checking. These
artifacts are sufficient to justify correctness of the execution, but they are
not succinct enough for direct blockchain verification. The purpose of the final
wrapper is therefore to compress them into a single proof $\pi_{final}$.

\subsection{Wrapped statement}

The wrapper proves an NP statement that can be denoted by a relation
$R_{wrap}(x_{pub}, w_{wrap}) = 1$. Its public input contains the program
commitment, the initial state commitment $z_{0\_\text{comm}}$, the claimed
final state commitment $z_{n\_\text{comm}}$, and the claimed public result of
the computation. The private witness contains:

\begin{itemize}
    \item the final accumulator instance-witness pair $(U_n, W_n)$,
    \item the final folding proof $\pi_{n\_\text{fold}}$,
    \item the auxiliary witness needed to re-check the last IVC transition,
    \item the deferred memory witnesses, including the original transcript
          $tr$ and its sorted form $tr'$.
\end{itemize}

The relation $R_{wrap}$ is satisfied if and only if all of the following hold:

\begin{enumerate}
    \item the final folded CCS accumulator is valid,
    \item the last folding proof is valid and consistent with that accumulator,
    \item the claimed final state commitment matches the state obtained after the
          last VM transition,
    \item the memory transcripts satisfy permutation, $(a,t)$-ordering, and
          read-over-write semantics.
\end{enumerate}

In other words, the wrapper turns the full end-of-execution claim into a single
algebraic satisfiability statement. The outer verifier no longer needs to inspect
the large accumulator or the complete memory log directly.

\subsection{SuperSpartan layer}

To prove satisfiability of $R_{wrap}$, the proposed design uses a
SuperSpartan-style IOP \cite{Spartan, HyperNova, CCS}. This is a natural fit,
because the internal objects of the zkVM are already expressed as CCS relations.
Instead of embedding yet another full SNARK verifier into the wrapper circuit,
SuperSpartan reduces the final CCS statement to a small number of sum-check and
polynomial evaluation claims.

Concretely, the witness of the wrapper contains the same algebraic objects that
appear at the end of the IVC process: accumulator coordinates, folding witness
variables, and memory-check witnesses. SuperSpartan verifies these through a
sum-check based reduction, leaving only polynomial evaluation queries to be
proven succinctly. This makes the wrapper conceptually aligned with the rest of
the design: the step circuit proves local execution, the IVC chain proves global
consistency, and the wrapper only compresses the resulting final relation.

\subsection{Greyhound layer}

The remaining polynomial evaluation claims are proven with Greyhound
\cite{Greyhound}. This choice is important for two reasons. First, Greyhound is
a lattice-based polynomial commitment scheme, so it preserves the post-quantum
assumption stack of the overall zkVM. Second, it provides succinct evaluation
proofs with sublinear verification complexity, which is exactly the final
compression step needed after SuperSpartan.

This means that the wrapper does not switch back to a discrete-log based or
trusted-setup polynomial commitment scheme. Instead, the final proof remains in
the same general lattice-based design space as the folding layer itself. The
Greyhound paper reports evaluation proofs around 53 KB for large polynomial
degrees, which suggests that a future fully implemented wrapper could plausibly
produce a final proof in the range of tens of kilobytes, rather than exposing
the much larger native accumulator.

\subsection{Composed soundness intuition}

The soundness of the final proof is obtained by composition. If
$\pi_{final}$ verifies, then the soundness of the SuperSpartan and Greyhound
wrapper implies that there exists a witness satisfying $R_{wrap}$. By
construction of $R_{wrap}$, this witness contains a valid final folded
accumulator and a valid final folding proof. By soundness of the LatticeFold
recursion, this implies the existence of a consistent unfolded IVC chain. By the
step circuit design from Section~\ref{sec:riscv_circuit}, each step in that
chain corresponds to a valid RISC-V transition. Finally, the deferred memory
proof ensures that all logged memory accesses satisfy read-over-write semantics.

Therefore, an accepting final wrapper proof implies existence of a valid
execution trace of the committed program, starting from the public initial state
and ending in the claimed final state.

\subsection{Status in this work}

The current prototype does not yet implement this final wrapper. The section
should therefore be understood as the design of the final compression layer, not
as an experimentally evaluated implementation artifact. In the intended
non-interactive instantiation, verifier challenges are derived using the
Fiat-Shamir heuristic, so the wrapper follows the standard random-oracle style
interpretation used by modern folding and SNARK systems.
